\section{Introduction}
\label{sec:introduction}

Memory updates are often compared by naming architectures or counting
parameters.  A transaction, however, imposes a more specific requirement: the
controller must realize a state change from the family demanded by the
environment.  A model can have ample parameters yet retain irreducible error
when its reachable update set lacks an erase/write magnitude, value channel,
address, basis, rank, or state-dependent freedom.

We therefore study the thesis that \emph{the geometry and algebra of update
demands determine the minimal memory-control architecture}.  Our experiments
separate reachable-set constraints from optimization, then add controller
freedoms one at a time.  The resulting evidence chain links behavioral
reachability, learned control rank, joint diagonalizability, a nested
architecture-demand lattice, and repeated structured transactions.

The paper makes six controlled contributions:
\begin{enumerate}
  \item a behavioral-reachability account of unseen update error;
  \item prospective learned-rank and joint-diagonalization tests of controller
        sufficiency;
  \item a static architecture-demand lattice spanning magnitude, value,
        address, and state-conditioned control;
  \item a repeated-sequence lattice with registered simpler-demand, retention,
        and long-gap direction guardrails;
  \item a paired tied-versus-dual sequence bridge and a fixed-codebook
        localization/candidate decomposition; and
  \item immutable artifact provenance and deterministic figure regeneration.
\end{enumerate}

These are controlled-reference and structured relational results.  They are
not passed official GDN2/KDA or KVEraser evidence, pretrained-language-model
evidence, or general agent evidence.  The systems measurement is reported
separately as an in-process controlled proxy.
